\section{New Years Message}

The ancient Tantra teachings say that the way to realization differs from age to age, so that methods appropriate to the golden age may lose their effectiveness in our crepuscular times, the ``Kali Yuga''. While passive contemplation was appropriate to a more spiritual age, we are weighed down with matter, our bodies are heavy, the mundane affairs of the world distract us. We can use that against itself and follow a path of action, practical realization, and direct experience.

Let's keep in mind that our true Self, the Atman, is more than the mere passive observer of appearances. It is their creator and their source, as they perpetually arise from that unknown nothingness, what we've been calling by that ugly word ``noumenon''. This is not action in the sense of frenetic activity, but action in the true sense, action as wei wu wei, action as ``letting it happen''.

We come to see that between thought and action there is no separation. Thought is possibility and action is thought made actual. Any other thought is not ``ours'' and no value must be attached to it.

A stone dropped from a height falls to the ground -- it cannot act against its nature. An animal may choose the means but its ends are determined by its nature. The human is otherwise. He can be in out of harmony with the Tao; he can neglect his dharma; he can reject the Logos.

As we come to realize our non-dual nature, we see that to speak of the human ``being'' is inadequate. For our ``being'' is incomplete without taking into account non-being ... the human state is always in potential. As our will actualizes itself toward its ideal, it simultaneously expresses, creates, and reveals our personality. That is how we come to know who we are.

Our personality is the heart of our being. As manifested, it is an ``individual'', one physical being among others. It becomes our error to then identify with our individuality in matter rather than our personality in spirit. We see ourselves as individuals vying with other individuals. This is clearly false because the solitary individual could not even survive and would not even have language for thought or culture, nor means for self-realization without a community of selves, both synchronically and diachronically.

We then know ourselves as a spiritual community each with our own possibilities to actualize. Ultimately, we are bringing the possibilities of the One into manifestation as the drama of creation. Some want to make this totally arbitrary as though they as individuals possessed more wisdom than the Logos through which all things are created. Others wonder why our possibilities are so different; but that is simply to ask why one being is not another, a nonsensical question. Others call it an injustice and would prefer total egalitarianism. But this is to deny the Infinity and Goodness of the One.

As we make our plans and resolutions for the New Year, keep in mind that because we act in a community of persons, that means the world is fundamentally a moral world as we create it through our collective choices. Are our choices free, arising from our deepest Self, or are they the playthings of external forces? Are we seeking to actualize our highest self or merely satisfying the inordinate desires of ourselves as individuals? During the ``twilight of the gods'', as we grope in darkness, sometimes struggling with death, disease, and disappointment, we need only turn back toward the Light.


\flrightit{Posted on 2008-01-01 by Cologero}

